\begin{afterword}
\summaryhead

The introduction to the fifth book, \textsc{Mere Goodness}, defines value theory as the study of what people care about, including what we wish we cared about and everyday values that are not moral. We study our values because we do not act as we wish, are not consistent enough to have utility functions, and misjudge even what we wish. Philosophers disagree about what we want, what we ought to want and what ``ought'' means, and the record of moral theory is poor.

The book's three sequences cover failed theories of value, obstacles to a new one, and how to apply such theories. The last is called the most important, since a normative theory is worth what it does for practice; when to trust intuition over theory is decided case by case.

Two questions will recur. Where should we go? Here the introduction presents Yudkowsky's ``fun theory,'' transhumanism, and the neglect of utopia-planning. And how do ends relate to means? It closes by asking what we build into the word ``valuable.''

\responsehead

The introduction previews the book and poses its questions, and most of what it says about positions and people checks out. The three sequences are the book's. Fun theory is Yudkowsky's term, used as early as 2002. The definition of transhumanism is close to the one in the Transhumanist FAQ. That human choices do not fit a consistent utility function is asserted without evidence here, but the book supplies it later, in the posts on the Allais paradox. The disagreement it reports among philosophers is borne out by the 2009 PhilPapers survey for moral theory, where none of the three main theories drew more than about a quarter of respondents; on metaethics a majority leaned toward moral realism. Its answer to when to trust intuition, that it depends on the case, matches the last sequence, which argues on both sides.

Two of its judgments are stated as premises. The first is that ``The cash value of a normative theory is how well it translates into normative practice,'' which is why the sequence on application is called the most important. That is close to the pragmatist view of theories, and uses the phrase William James used for it; it is given no argument. The second is that utopia-planning is neglected, which has no source. The paragraph that follows argues, with a hedge, against giving up on utopia. It does not name a well-known case on the other side, Popper's preference for piecemeal reform over a plan for the good society. An introduction need not argue the other side, but the reader should know there is one.

A few terms are used in the book's own sense. ``Value theory'' is defined by what people care about, where philosophers use it for the theory of what is good. ``Hedonism'' and ``eudaimonia'' are glossed briefly. In the standard classification of theories of well-being, the rivals to hedonism are desire theories and objective-list theories, and eudaimonist accounts are counted among the latter. That is useful context for the next post, which argues against hedonism. Sandel's objection to enhancement, glossed as making life feel less of a ``gift,'' is in Sandel's words ``the hubris objection,'' about the disposition that enhancement expresses.

In short: an introduction whose descriptions of positions mostly check out, whose ranking of practice above theory rests on a pragmatist premise it does not argue, and whose claim that utopia-planning is neglected has no source.
\end{afterword}
